\documentclass[a4,11pt]{aleph-notas}

% -- Paquetes adicionales
\usepackage{enumitem}
\usepackage{aleph-comandos}

% -- Datos
\institucion{Facultad de Ciencias Exactas, Naturales y Ambientales}
\carrera{Catálogo STEM}
\asignatura{Matemática III}
\tema{Resumen 12: Cambio de variable}
\autor[A. Merino]{Andrés Merino}
\fecha{Periodo 2026-1}

\logouno[0.16\textwidth]{Logos/logoPUCE_04_ac}
\definecolor{colortext}{HTML}{1957d1}
\definecolor{colordef}{HTML}{1957d1}
\fuente{montserrat}


\begin{document}

\encabezado

%%%%%%%%%%%%%%%%%%%%%%%%%%%%%%%%%%%%%%
\section{Cambio de variable en integrales dobles}
%%%%%%%%%%%%%%%%%%%%%%%%%%%%%%%%%%%%%%

\begin{advertencia}
    Sea
    \[
        \funcion{T}{\R^2}{\R^2}{(u,v)}{T(u,v)}
    \]
    una aplicación lineal con representante
    \[
        [T]= A =\begin{pmatrix}
            a & b \\
            c & d
          \end{pmatrix},
    \]
    tal que $\det(A)\neq 0$. Entonces la función es biyectiva, envía paralelogramos en paralelogramos y vértices de paralelogramos en vértices de paralelogramos. Además, si $\Omega^*$ es un paralelogramo, entonces $\Omega=T(\Omega^*)$ es un paralelogramo y
    \[
        \text{Área de } \Omega = |\det(A)| \cdot (\text{Área de } \Omega^*).
    \]
\end{advertencia}

En adelante, tomaremos la función inyectiva y diferenciable con continuidad
\[
    \funcion{T}{\R^2}{\R^2}{(u,v)}{T(u,v)=(x(u,v),y(u,v))}
\]
a la cual llamaremos una transformación de coordenadas.

\begin{prop}
    Sean $\Omega$ y $\Omega^*$ regiones elementales de $\R^2$. Supongamos que $\func{T}{\R^2}{\R^2}$ es una transformación de coordenadas tal que $T(\Omega^*)=\Omega$. Si $\func{f}{\Omega}{\R}$ es una función integrable, entonces
    \[
        \iint_\Omega f(x,y)\, dxdy = \iint_{\Omega^*} f(T(u,v)) |\det(J_T(u,v))| \, dudv.
    \]
\end{prop}

\begin{advertencia}
    Si
    \[
        \funcion{T}{\R^2}{\R^2}{(u,v)}{T(u,v)=(x(u,v),y(u,v))}
    \]
    es una transformación de coordenadas biyectiva, entonces
    \[
        \funcion{T^{-1}}{\R^2}{\R^2}{(x,y)}{T^{-1}(x,y)=(u(x,y),v(x,y))}
    \]
    también es una transformación de coordenadas; además,
    \[
        \det(J_{T^{-1}}(x,y)) = \frac{1}{\det(J_{T}(u,v))},
    \]
    con $(x,y)=T(u,v)$.
\end{advertencia}

%%%%%%%%%%%%%%%%%%%%%%%%%%%%%%%%%%%%%%
\section{Coordenadas polares}
%%%%%%%%%%%%%%%%%%%%%%%%%%%%%%%%%%%%%%

La transformación de coordenadas polares es
\[
    \funcion{T}{\R^2}{\R^2}{(r,\theta)}{(r\cos\theta,\, r\sin\theta)},
\]
con $r\geq 0$ y $\theta\in[0,2\pi)$. Su jacobiano es
\[
    \det(J_T(r,\theta)) = r.
\]

Por tanto, si $\Omega^*$ es la región en coordenadas polares correspondiente a $\Omega$, entonces
\[
    \iint_\Omega f(x,y)\, dxdy = \iint_{\Omega^*} f(r\cos\theta,\, r\sin\theta)\, r\, drd\theta.
\]

\begin{advertencia}
    Las regiones circulares o sectoriales se describen de forma natural en coordenadas polares. Algunos ejemplos:
    \begin{itemize}
        \item Disco de radio $a$: $\Omega^*=\{(r,\theta):0\leq r\leq a,\, 0\leq\theta\leq 2\pi\}$.
        \item Sector circular: $\Omega^*=\{(r,\theta):0\leq r\leq a,\, \alpha\leq\theta\leq\beta\}$.
        \item Corona circular: $\Omega^*=\{(r,\theta):a\leq r\leq b,\, 0\leq\theta\leq 2\pi\}$.
    \end{itemize}
\end{advertencia}


\end{document}
